
%=====================================================================
\section{$L(0)$ identities for the cusp rows}\label{sec:L0}
%=====================================================================
For a cusp row $c$, $\kappa_c(x)$ is the $t^0$ coefficient of $f_c(\e(x)e^{-t})$; by Mellin, $\kappa_c(x)=Z_c(0;x)$ with $Z_c(w;x)=\sum A_Q\e(Qx)Q^{-w}$ entire.
For $x=h/k$ with $k$ composed of $2,5,7$, $\e(h\Nm\fa/(140k))$ is a ray class function mod $\fm_k$ (checked on all classes to norm $60\,000$); for $3\mid k$ the
ideals divisible by the inert prime $3$ give one extra ray-class sum. This yields
\[
\kappa_c(h/k)=\sum_{\chi'}\widehat G_{c,h/k}(\chi')\,L_{\fm_k}(0,\chi'),\]
and, for $3\nmid k$,
\[
\kappa_c(h/k)=\frac1{\varphi(140k)}\sum_{\psi\bmod140k}\tau(\bar\psi)\psi(h)\sum_\chi\widehat F_c(\chi)\,L(0,u_\chi\otimes\psi).
\]
Only totally odd characters contribute; for those $L(0,\chi)=W(\chi)\cdot140\cdot L(1,\bar\chi)/\pi^2$ (checked to $2\cdot10^{-15}$). Type-$(0,0)$ and mixed-type
characters have $L(0)=0$ (Category A from the exact PARI-free Shintani sums of Appendix~\ref{app:pari}; the PARI values agree).
\catB{Numerically verified, conditional on PARI \texttt{lfun} (38 digits; primitive characters via \texttt{bnrmap}, imprimitive Euler factors restored): all
$24$ cusp rows at all $7$ points to $4\cdot10^{-14}$ (Table~\ref{tab:phi73}). Not proved. Central-value hypothesis rejected: replacing $L(0)$ by $L(\tfrac12)$
gives $\Phi_{73}(1)=0.0349+0.0748i$, max deviation $0.32$ ($0.31$ with only the $12$ totally odd central values). At $x=1$ the $24$ values lie in a
$6$-dimensional space.}

\paragraph{Hand certification of $\mathrm{Cl}^+_{(2\sqrt{35})\infty_1\infty_2}$.} From
$1\to U/U_{\fm,1}\to(\Ok/\fm)^\times\times\{\pm\}^2\to\mathrm{Cl}^+_\fm\to\mathrm{Cl}\to1$ with $(\Ok/(2))^\times=\langle\sqrt{35}\rangle\cong\Z/2$,
$\mathbb F_5^\times=\langle2\rangle$, $\mathbb F_7^\times=\langle3\rangle$, $U=\langle-1,\varepsilon\rangle$, $\mathrm{Cl}=\Z/2$ generated by $\fp_{13}$
($\fp_{13}^2=(22+3\sqrt{35})$), the Smith normal form gives invariants $[2,2,24]$, order $96=2\cdot48\cdot4/4$. The $\sigma$-action, infinity types and conductors
are read off and match PARI. (Category A, PARI-free.)

\begin{table}[ht]\centering\small
\begin{tabular}{llll}\toprule
$\chi$ ($\sigma$-pair) & type & $L(\tfrac12,u_\chi)$ & $L(0)$\\\midrule
$(1,\cdot,1)\sim(13,\cdot,1)$ & $(0,0)$ & $1.109240202417-0.205788403698i$ & $0$\\
 & $(1,1)$ & $0.131513747745+0.708885114750i$ & $-(6+4\sqrt3)+(6\sqrt3-4)i$\\
$(3,\cdot,1)\sim(15,\cdot,1)$ & $(0,0)$ & $1.057777176883-0.300492329005i$ & $0$\\
 & $(1,1)$ & $-0.123002454086-0.432986722370i$ & $-12+4i$\\
$(5,\cdot,1)\sim(17,\cdot,1)$ & $(0,0)$ & $0.514021161380+0.466537408032i$ & $0$\\
 & $(1,1)$ & $1.122786592523-1.237062791397i$ & $(4\sqrt3-6)-(4+6\sqrt3)i$\\\bottomrule
\end{tabular}
\caption{Central values and $L(0)$ for six of the twelve dihedral forms; the other six ($j=7,9,11\sim19,21,23$) are complex conjugates.
All $12$ central values are nonzero. PARI \texttt{lfun}, 38 digits; $L(0)$ values recognised to $10^{-10}$.}\label{tab:L}
\end{table}

\subsection{PARI-free certification}\label{app:pari}
\catA{An exact, finite computation, or a standard theorem applied to exact data. Hand inputs: $\varepsilon=6+\sqrt{35}$ is the
fundamental unit and $\mathrm{Cl}=\Z/2$ (generated by $\fp_{13}$, $\fp_{13}^2=(22+3\sqrt{35})$).}
\begin{itemize}[leftmargin=2em]
\item The ray class group mod $\fm=\fp_2^3\fp_5\fp_7\infty_1\infty_2$, constructed by hand from
$U\to(\Ok/\fm)^\times\times\{\pm\}^2\to\mathrm{Cl}^+_\fm\to\mathrm{Cl}\to1$, has order $96$ and structure $\Z/24\times\Z/2\times\Z/2$; the same
holds for $\ff=(2\sqrt{35})\infty_1\infty_2$.
\item The independent class assignment of all $140$ components reproduces the PARI fit exactly: the $22$ empty (label, sign) slots, every
content $g$, the Eisenstein/cusp/mixed norms (differences $\le 6.7\cdot10^{-16}$), the same $24$ cusp components with coefficients
$\tfrac1{48}\e(k/60)$, and the label-row ranks $60$ and $12$ (in function space $36$ and $6$, \S\ref{sec:spec}). The theta identity holds for all $n$ with $|Q|\le400$ (residual
$\le8.9\cdot10^{-16}$); the statement for all $n$ is Category B.
\item Exact Shintani cone sums give the $12$ values $L(0,\chi)$ of the totally odd, non-$\sigma$-invariant characters of conductor
$(2\sqrt{35})$ in $\Z[\zeta_{12}]$: $-(6+4\sqrt3)\pm(6\sqrt3-4)i$, $-12\pm4i$, $(4\sqrt3-6)\mp(4+6\sqrt3)i$, each twice (one per $\sigma$-pair).
Validations: $\sum_c\zeta(0,c)=0$; $L(0,\chi)=0$ for all $72$ non-totally-odd characters; $L(\chi)=L(\chi\circ\sigma)$; a genus test exact for all
$8$ characters mod $20$; two independent moduli agree exactly.
\item All $7\times24$ cusp-row identities of Theorem~\ref{thm:K} hold to $1.2\cdot10^{-14}$ with PARI-free inputs. No discrepancy was found.
\end{itemize}
\emph{Not} redone without PARI: $L(\tfrac12,\chi)$, $L(1,\chi)$, root numbers. Characters are determined only up to $\chi\leftrightarrow\chi\circ\sigma$.
Note that the group mod $\fm$ (order $96$) is not the group mod $\fp_2\fp_5\infty_1\infty_2$ (order $16$, $\Z/8\times\Z/2$) used in \S\ref{sec:rq}.
